% Faithful formal-publication context for search/count reductions: Sections1.2 and2.
\section*{Author notation and isometry models}
\paragraph{Vector spaces.} Let $\mathbb F$ be a field. In this paper we only consider finite-dimensional vector spaces over $\mathbb F$. We use $\mathbb F^n$ to denote the vector space of length-$n$ column vectors. The $i$th standard basis vector of $\mathbb F^n$ is denoted $\vec e_i$. Depending on the context, 0 may denote the zero vector space, a zero vector, or an all-zero matrix. For $S$ a set of vectors, we use $\langle S\rangle$ to denote the subspace spanned by elements in $S$.
\paragraph{Some groups.} The general linear group of degree $n$ over a field $\mathbb F$ is denoted by $\mathrm{GL}(n,\mathbb F)$. The symmetric group of degree $n$ is denoted by $S_n$. The natural embedding of $S_n$ into $\mathrm{GL}(n,\mathbb F)$ is to represent permutations by permutation matrices. The subgroup of $\mathrm{GL}(n,\mathbb F)$ consisting of diagonal matrices is called the diagonal subgroup, denoted by $\operatorname{diag}(n,\mathbb F)$. A monomial matrix is a product of a diagonal and a permutation matrix; equivalently, each row and each column has exactly one non-zero entry. The collection of monomial matrices forms a subgroup of $\mathrm{GL}(n,\mathbb F)$, which we call the monomial subgroup and denote by $\operatorname{Mon}(n,\mathbb F)$. It is the semi-direct product $\operatorname{diag}(n,\mathbb F)\rtimes S_n\cong(\mathbb F^*)^n\rtimes S_n$.
\paragraph{Matrices.} Let $\mathrm M(\ell\times n,\mathbb F)$ be the linear space of $\ell\times n$ matrices over $\mathbb F$, and $\mathrm M(n,\mathbb F):=\mathrm M(n\times n,\mathbb F)$. Given $A\in\mathrm M(\ell\times n,\mathbb F)$, $A^t$ denotes the transpose of $A$.
A matrix $A\in\mathrm M(n,\mathbb F)$ is alternating, if for any $u\in\mathbb F^n$, $u^tAu=0$. That is, $A$ represents an alternating bilinear form. Note that in characteristic $\ne2$, alternating is the same as skew-symmetric, but in characteristic 2 they differ (in characteristic 2, skew-symmetric=symmetric). The linear space of $n\times n$ alternating matrices over $\mathbb F$ is denoted by $\Lambda(n,\mathbb F)$.
The $n\times n$ identity matrix is denoted by $I_n$, and when $n$ is clear from the context, we may just write $I$. The elementary matrix $E_{i,j}$ is the matrix with the $(i,j)$th entry being 1, and other entries being 0. The $(i,j)$-th elementary alternating matrix is the matrix $E_{i,j}-E_{j,i}$.
\paragraph{Matrix tuples.} We use $\mathrm M(\ell\times n,\mathbb F)^m$ to denote the linear space of $m$-tuples of $\ell\times n$ matrices. Boldface letters like $\mathbf A$ and $\mathbf B$ denote matrix tuples. Let $\mathbf A=(A_1,\ldots,A_m)$, $\mathbf B=(B_1,\ldots,B_m)\in\mathrm M(\ell\times n,\mathbb F)^m$. Given $P\in\mathrm M(\ell,\mathbb F)$ and $Q\in\mathrm M(n,\mathbb F)$, $P\mathbf A Q:=(PA_1Q,\ldots,PA_mQ)\in\mathrm M(\ell,\mathbb F)$. Given $R=(r_{i,j})_{i,j\in[m]}\in\mathrm M(m,\mathbb F)$, $\mathbf A^R:=(A'_1,\ldots,A'_m)\in\mathrm M(m,\mathbb F)$ where $A'_i=\sum_{j\in[m]}r_{j,i}A_j$.
\begin{authorclaim}{Remark5}
In particular, note that the coefficients in the formula of defining $A'_i$ correspond to the entries in the $i$th column of $R$. While this choice is immaterial (we could have chosen the opposite convention), all of our later calculations are consistent with this convention.
\end{authorclaim}
Given $\mathbf A,\mathbf B\in\mathrm M(\ell\times n,\mathbb F)^m$, we say that $\mathbf A$ and $\mathbf B$ are isometric, if there exists $P\in\mathrm{GL}(n,\mathbb F)$, such that $P^t\mathbf A P=\mathbf B$. Finally, $\mathbf A$ and $\mathbf B$ are pseudo-isometric if there exist $P\in\mathrm{GL}(n,\mathbb F)$ and $R\in\mathrm{GL}(m,\mathbb F)$, such that $P^t\mathbf A P=\mathbf B^R$.
\paragraph{Matrix spaces.} Linear subspaces of $\mathrm M(\ell\times n,\mathbb F)$ are called matrix spaces. Calligraphic letters like $\mathcal A$ and $\mathcal B$ denote matrix spaces. By a slight abuse of notation, for $\mathbf A\in\mathrm M(\ell\times n,\mathbb F)^m$, we use $\langle\mathbf A\rangle$ to denote the subspace spanned by those matrices in $\mathbf A$. For $\mathbf A,\mathbf B\in\mathrm M(n,\mathbb F)^m$, we say that the spaces $\langle\mathbf A\rangle,\langle\mathbf B\rangle$ are isometric iff the tuples $\mathbf A,\mathbf B$ are pseudo-isometric.
\begin{authorclaim}{Definition4 (The Alternating Matrix Space Isometry problem)}
Input: $A_1,\ldots,A_m$ and $B_1,\ldots,B_m$, $n\times n$ alternating matrices over a field $\mathbb F$.
Decide: Is there a $L\in\mathrm{GL}(n,\mathbb F)$, such that the linear span of $\{A_i:i\in[m]\}$ is equal to the linear span of $\{L^tB_iL:i\in[m]\}$?
\end{authorclaim}
\paragraph{3-way arrays.} Let $\mathrm T(\ell\times n\times m,\mathbb F)$ be the linear space of $\ell\times n\times m$ 3-way arrays over $\mathbb F$. We use the fixed-width teletypefont for 3-way arrays, like $\mathtt A,\mathtt B$, etc..
Given $\mathtt A\in\mathrm T(\ell\times n\times m,\mathbb F)$, we can think of $\mathtt A$ as a 3-dimensional table, where the $(i,j,k)$th entry is denoted as $\mathtt A(i,j,k)\in\mathbb F$. We can slice $\mathtt A$ along one direction and obtain several matrices, which are then called slices. For example, slicing along the first coordinate, we obtain the horizontal slices, namely $\ell$ matrices $A_1,\ldots,A_\ell\in\mathrm M(n\times m,\mathbb F)$, where $A_i(j,k)=\mathtt A(i,j,k)$. Similarly, we also obtain the lateral slices by slicing along the second coordinate, and the frontal slices by slicing along the third coordinate.
We will often represent a 3-way array as a matrix whose entries are vectors. That is, given $\mathtt A\in\mathrm T(\ell\times n\times m,\mathbb F)$, we can write
\[\mathtt A=\begin{bmatrix}w_{1,1}&w_{1,2}&\cdots&w_{1,n}\\w_{2,1}&w_{2,2}&\cdots&w_{2,n}\\\vdots&\vdots&\ddots&\vdots\\w_{\ell,1}&w_{\ell,2}&\cdots&w_{\ell,n}\end{bmatrix},\]
where $w_{i,j}\in\mathbb F^m$, so that $w_{i,j}(k)=\mathtt A(i,j,k)$. Note that, while $w_{i,j}\in\mathbb F^m$ are column vectors, in the above representation of $\mathtt A$, we should think of them as along the direction ``orthogonal to the paper.'' Following [45], we call $w_{i,j}$ the tube fibers of $\mathtt A$. Similarly, we can have the row fibers $v_{i,k}\in\mathbb F^n$ such that $v_{i,k}(j)=\mathtt A(i,j,k)$, and the column fibers $u_{j,k}\in\mathbb F^\ell$ such that $u_{j,k}(i)=\mathtt A(i,j,k)$.
Given $P\in\mathrm M(\ell,\mathbb F)$ and $Q\in\mathrm M(n,\mathbb F)$, let $P\mathtt A Q$ be the $\ell\times n\times m$ 3-way array whose $k$th frontal slice is $PA_kQ$. For $R=(r_{i,j})\in\mathrm{GL}(m,\mathbb F)$, let $\mathtt A^R$ be the $\ell\times n\times m$ 3-way array whose $k$th frontal slice is $\sum_{k'\in[m]}r_{k',k}A_{k'}$. Note that these notations are consistent with the notations for matrix tuples above, when we consider the matrix tuple $\mathbf A=(A_1,\ldots,A_k)$ of frontal slices of $\mathtt A$.
